\section*{Abstract}
Botnets pose a significant cybersecurity threat, enabling attacks such as DDoS, data theft, and service disruptions on IoT devices. These devices often lack built-in botnet traffic filtering, leaving them highly exposed. Existing AI-based solutions improve detection capabilities but have limitations: (i) they are too heavy for IoT deployment, and (ii) they lack unlearning capabilities to forget sensitive or outdated features without retraining. To address these challenges, we propose DiRLU, a lightweight, reinforcement learning driven framework, while ensuring privacy by selectively unlearning sensitive or outdated features without requiring retraining. The framework leverages knowledge distillation to transfer knowledge from a teacher model into a lightweight student model, with both models trained using A2C. A post-hoc unlearning mechanism modifies weights to remove targeted features, while restored features show negligible performance loss, confirming reversibility. Unlike many benchmark models that used only 5\% of the BoT-IoT dataset, this research leverages 25\%, allowing us to develop a strong teacher model. Both the teacher and student models were trained using the A2C reinforcement learning algorithm, achieving impressive results, with the student model achieving 99.60\% accuracy and a 99.80\% F1 score. To enhance transparency, we integrated Explainable AI (XAI), particularly LIME, which helps interpret the model’s decisions and identify the key features influencing its predictions. DiRLU combines efficiency with privacy, aligning with GDPR standards (right to be forgotten) to provide practical IoT security. By combining knowledge distillation, feature unlearning and XAI, this research not only strengthens botnet detection but also sets new standards for security, interpretability, and data privacy in cybersecurity.

\vspace{1cm}
\textbf{Keywords:} Botnet, Reinforcement Learning, Feature Unlearning, Knowledge Distillation, Explainable AI, A2C, Security \& Privacy
\pagebreak
